\appendix
\section{שחזור והרצה}

המחברת והדוח בנויים כך שאפשר לעיין בתוצאות המאומתות בלי לאמן מחדש את כל המודלים. התוצאות המספריות והאיורים נטענים מקבצי תוצרים שנשמרו מהריצות הסופיות.

שתי הריצות המרכזיות הן:

\begin{itemize}
    \item ריצת HMM מורחבת: \code{experiments_extended/extended_20260811_121957/}
    \item ריצת מודלי \tech{Supervised Learning}: \code{experiments_supervised/supervised_20260811_133344/}
\end{itemize}

סביבת הריצה המאומתת השתמשה ב־\tech{Python 3.11.15}. בריצת ה־HMM נעשה שימוש ב־\tech{hmmlearn 0.3.3}, \tech{NumPy 2.4.6}, \tech{pandas 3.0.3} ו־\tech{scikit-learn 1.9.0}.

להרצה מחדש של הניסוי המלא ניתן להשתמש בפקודות הבאות:

\begin{english}
\begin{verbatim}
python -m pip install -r requirements.txt
python run_extended_analysis.py
python analyze_extended_context.py \
  --run-dir experiments_extended/<run_id>
python run_supervised_baselines.py
jupyter nbconvert --to notebook --execute \
  HMM_Market_Regimes_Project.ipynb \
  --output /tmp/HMM_verified.ipynb \
  --ExecutePreprocessor.timeout=600
\end{verbatim}
\end{english}

כל ריצה חדשה נכתבת לתיקייה חדשה בעלת חותמת זמן ואינה דורסת את התוצרים שעליהם מבוסס הדוח. ההפרדה הזו חשובה משום שמקור הנתונים חיצוני, ונתוני Yahoo Finance או זמינות ההורדה יכולים להשתנות לאורך זמן. כך אפשר להבדיל בין התוצאות ששימשו להגשה לבין ריצה עתידית חדשה.
